\documentclass[10pt,a4paper]{amsart}
\usepackage[margin=19mm]{geometry}
\usepackage{amsmath,amssymb,mathtools}
\usepackage[scr=rsfs,cal=euler]{mathalfa}
\usepackage{xfrac}
\usepackage[unicode,hidelinks,bookmarks=false]{hyperref}
\newcommand{\ie}{i.e.\ }
\newcommand{\cf}{cf.\ }
\newcommand{\resp}{respectively\ }
\newcommand{\Subsection}{\subsection}
\providecommand{\nobreakdash}{\nobreak\hbox{-}}
\setlength{\emergencystretch}{2em}
\multlinegap0pt
\usepackage{bm}
\usepackage{bbm}
\usepackage{dsfont}
\usepackage{xspace}
\usepackage{cancel}
\usepackage{mathtools}
\usepackage{amsthm}
\makeatletter
\@ifundefined{theorem}{\newtheorem{theorem}{Theorem}}{}
\@ifundefined{lemma}{\newtheorem{lemma}[theorem]{Lemma}}{}
\makeatother
\newcommand{\CZ}{\mathbb{C}\mathfrak{Z}}
\newcommand{\Mat}{{\rm Mat}}
\newcommand{\mI}{\mathbb{I}}
\title[A Spectral Theorem for Zeon Matrices]{A Spectral Theorem for Zeon Matrices}
\author{G. Stacey Staples}
\date{Source: arXiv:2201.09321v3 (CC BY 4.0)}
\begin{document}
\maketitle

\noindent\emph{Source and licence.} A Spectral Theorem for Zeon Matrices. Version arXiv:2201.09321v3 (CC BY 4.0), distributed under the Creative Commons Attribution 4.0 International licence, as recorded in the arXiv licence metadata for this exact version and shown on the abstract page https://arxiv.org/abs/2201.09321v3. Full rights notice: licenses/arxiv-2201.09321-CC-BY-4.0.txt.\par\smallskip

\noindent\textit{Editorial scope.} Counted unit: the Theorem whose source label is \texttt{row reduction II} in the article (source0.tex lines 654--656, source label \texttt{row reduction II}), delivered together with the article's own complete original proof of it (source0.tex lines 658--678, one \texttt{proof} environment of 21 lines). No mathematics is written, repaired or re-derived: statement and proof are the article's own text line for line. Required context retained uncounted: row reduction I (lines 619--636). Every definition, display and label of those lines is retained unchanged. Omitted from this package and not redistributed: 1--618, 637--653, 657--657, 679--1349. Nothing inside a retained excerpt is omitted or altered: every retained body is delivered exactly as the article's lines are written, so no omission receipt of any kind accompanies this package. Citations: the retained excerpts carry no citation command at all, so no bibliography entry of the article is transcribed and no reference is redistributed. Reference adaptation: none was needed and none was made; every reference inside the delivered unit resolves inside it, so no unresolved reference or citation is produced. Printed slips kept literal: no printed word, symbol, hypothesis or number of the counted statement or of its proof was altered. Layout and class: the article's own class is \texttt{article} with packages including authblk, amsmath, amssymb, amsthm, graphicx, url, enumerate, algorithm2e; the wrapper is the lane's pinned \texttt{amsart} editorial wrapper at 10pt on A4, which restates the theorem-like environments the retained text needs and adds the article's own macro definitions that the retained text needs (\texttt{\textbackslash CZ}, \texttt{\textbackslash Mat}, \texttt{\textbackslash mI}), with extra wrapper packages (bm, bbm, dsfont, xspace, cancel, mathtools). The delivered numbering is the unit's own. Assets: no retained span contains a figure environment, a graphics-inclusion command, a PDF-page inclusion, a table or a TikZ picture; no image file is copied into this package, the package declares no asset dependency and no image file is redistributed with it. Licence: version arXiv:2201.09321v3 of this article is distributed under the Creative Commons Attribution 4.0 International licence, as recorded in the arXiv licence metadata for this exact version and shown on the abstract page https://arxiv.org/abs/2201.09321v3. A full rights notice is delivered at licenses/arxiv-2201.09321-CC-BY-4.0.txt. Characters: every retained excerpt is pure ASCII, so the delivered engine, XeLaTeX with its default Latin Modern text font, reports no missing character.
\begin{lemma}\label{row reduction I}
Let $A\in\Mat(m;\CZ)$ be a zeon matrix of rank $m-1$.  If the $k$th column of $A$ consists entirely of nilpotent zeon elements, then there exists an invertible zeon matrix $Q$ such that $|QA|=\pm |A|$ and $QA$ is of the form \begin{equation}\label{QA matrix}
QA=\begin{pmatrix}
\alpha_1&0 &\cdots&0&\eta_1&0&0&\cdots&0\\
0&\alpha_2 & \cdots&0&\eta_2&0&0&\cdots&0\\
0&0&\ddots&0&\vdots&0&0&\cdots&0\\
0&0&\cdots&\alpha_{k-1}&\eta_{k-1}&0&0&\cdots&0\\
0&0&\cdots&0&\eta_{k}&0&0&\cdots&0\\
0&0&\cdots&0&\eta_{k+1}&\alpha_{k+1}&0&\cdots&0\\
0&0&\cdots&0&\eta_{k+2}&0&\alpha_{k+2}&\cdots&0\\
0&0&\cdots&0&\vdots&0&0&\ddots&0\\
0&0&\cdots&0&\eta_{m}&0&\cdots&0&\alpha_m\\
\end{pmatrix},
\end{equation}
where $\eta_1, \ldots, \eta_m\in\CZ^\circ$ and $\alpha_\ell\in\CZ^\times$ for $\ell\ne k$.  In particular, the determinant of $QA$ is \begin{equation*}
|QA|=\eta_k\prod_{\ell\ne k}\alpha_\ell.
\end{equation*}
\end{lemma}

\begin{theorem}\label{row reduction II}
Let $A\in\Mat(m;\CZ)$.  If $\lambda\in\CZ$ is a spectrally simple zero of $\chi_A(u)$, then there exists $\xi\in{\CZ}^m\setminus ({\CZ}^\circ)^m$ such that  $A\xi=\lambda\xi$.
\end{theorem}

\begin{proof}
Since $\lambda$ is a spectrally simple zero of the characteristic polynomial $\chi_A(u)$, the matrix $\lambda\mI-A$ is of rank one and thus, by applying Theorem \ref{row reduction I} and multiplying rows by element inverses, can be placed into the form 
\begin{equation*}
\rho(\lambda\mI-A)=\begin{pmatrix}
1&0 &\cdots&0&{\eta_1}&0&0&\cdots&0\\
0&1& \cdots&0&{\eta_2}&0&0&\cdots&0\\
0&0&\ddots&0&\vdots&0&0&\cdots&0\\
0&0&\cdots&1&{\eta_{k-1}}&0&0&\cdots&0\\
0&0&\cdots&0&\eta_k&0&0&\cdots&0\\
0&0&\cdots&0&{\eta_{k+1}}&1&0&\cdots&0\\
0&0&\cdots&0&{\eta_{k+2}}&0&1&\cdots&0\\
0&0&\cdots&0&\vdots&0&0&\ddots&0\\
0&0&\cdots&0&{\eta_{m}}&0&\cdots&0&1
\end{pmatrix},
\end{equation*}
where $\eta_1, \ldots, \eta_m\in\CZ^\circ$.  In particular, $\eta_k=\vert\rho(\lambda\mI-A)\vert=0$.  Solving the equation $(\lambda\mI-A)v=0$ now gives \begin{equation*}
v=\alpha\left(-{\eta_1},\cdots,-{\eta_{k-1}},1,-{\eta_{k+1}},\cdots,{\eta_m}\right)^\intercal,
\end{equation*}
satisfying $Av=\lambda v$ for any $\alpha\in\CZ$.

\end{proof}

\end{document}
